% Бүлэг 1 — Онол, арга зүйн судалгаа
% Шинэчилсэн: 2026-10-05. Эх template: ~/Downloads/template/Chapters/Chapter1.tex
% Тэмдэглэгээ: [А] = Нармандах, [Б] = Мэнд-Амар, [Х] = хамтарсан

\chapter{Шаардлага болон хичээлийн материалын чанарын үнэлгээний онол, арга зүйн судалгаа}
\label{Chapter1}

%-------------------------------------------------------------------------------
\newcommand{\keyword}[1]{\textbf{#1}}
\newcommand{\tabhead}[1]{\textbf{#1}}
\newcommand{\code}[1]{\texttt{#1}}
\newcommand{\file}[1]{\texttt{\bfseries#1}}
\newcommand{\option}[1]{\texttt{\itshape#1}}
\newcommand{\ph}{\textit{[Ноорог бичигдэнэ]}}
%-------------------------------------------------------------------------------

\section{Системийн тухай} % [Х] Багшийн заавар: 1-р хэсэгт өөрийн системийн дэлгэрэнгүй
% Систем юу хийх, хэн ашиглах, оролт/гаралт (шаардлагын баримт, хичээлийн материал -> илрүүлэлт, тайлбар, матриц).
\ph

\section{Програм хангамжийн шаардлагын инженерчлэл} % [А]
\subsection{Шаардлагын ойлголт ба төрөл}
\ph
\subsection{Шаардлагын чанарын шинж чанарууд}
% ISO/IEC/IEEE 29148: нэг утгатай, шалгагдах, бүрэн, тууштай гэх мэт.
\ph
\subsection{Шаардлагын чанарын алдааны төрлүүд}
% Бүрхэг (хэмжих боломжгүй), давхардсан, зөрчилдсөн. Жишээ хүснэгттэй.
\ph
\subsection{Шаардлагын алдааны өртөг}
% Хожуу илрэх тусам засах зардал өсөх. [эх сурвалж хэрэгтэй]
\ph

\section{Хичээлийн материалын уялдааны онол} % [Б]
\subsection{Суралцахуйн зорилт ба түүний ангилал}
\ph
\subsection{Зорилт, сургалтын үйл ажиллагаа, үнэлгээний уялдаа}
\ph
\subsection{Хамрах хүрээний матриц}
\ph

\section{Байгалийн хэлний боловсруулалтын аргууд} % [Х]
\subsection{Дүрэмд суурилсан арга}
\ph
\subsection{Текстийн embedding ба утгын ижил төстэй байдал}
% LaBSE, multilingual-e5; косинус ижил төстэй байдал, босго.
\ph
\subsection{Transformer загвар ба fine-tune}
% XLM-RoBERTa.
\ph
\subsection{Их хэлний загвар}
% LLM (Large Language Model), prompt, JSON schema, temperature.
\ph
\subsection{Байгалийн хэлний дүгнэлт}
% NLI (Natural Language Inference) — зөрчил илрүүлэх нэмэлт хэсэгт.
\ph

\section{Монгол хэлний боловсруулалтын онцлог} % [Х]
\subsection{Монгол хэлний бүтэц ба үгийн хувилгавар}
\ph
\subsection{Бага нөөцтэй хэл ба олон хэлт загварууд}
\ph

\section{Ижил төстэй судалгаа ба систем} % [А]+[Б]
\subsection{Шаардлагын чанарыг автоматаар үнэлэх олон улсын судалгаа} % [А]
\ph
\subsection{Хичээлийн материалын уялдааг шалгах судалгаа} % [Б]
\ph
\subsection{Монгол хэл дээрх судалгаа ба өмнөх дипломын ажлууд} % [Х]
\ph
\subsection{Судлагдсан аргуудын харьцуулалт}
% Харьцуулсан хүснэгт (longtable) + тайлбар.
\ph

\section{Үнэлгээний хэмжүүр} % [Х]
\subsection{Precision, Recall, F1}
\ph
\subsection{Тэмдэглэгээний тохироо: Cohen's kappa}
\ph
\subsection{Тогтвортой байдал, хурд, зардлын хэмжүүр}
\ph

\section{Технологийн судалгаа} % [Х]
\subsection{Машин сургалт ба backend технологи}
% Python, pandas, Transformers, sentence-transformers, scikit-learn, PyTorch, FastAPI.
\ph
\subsection{Frontend технологи}
% React.
\ph
\subsection{Өгөгдлийн сан}
% PostgreSQL.
\ph
\subsection{Хөгжүүлэлтийн орчин ба кодын менежмент}
% Docker, GitHub (branch + pull request), GitHub Projects, pytest, Google Colab.
\ph

% Бүлгийн дүгнэлт — багшийн заавраар ДУГААРГҮЙ
\section*{Бүлгийн дүгнэлт}
\addcontentsline{toc}{section}{Бүлгийн дүгнэлт}
\textit{[Бүлэг дууссаны дараа бичигдэнэ]}
